\chapter{Preface}

\section*{Why Would Anyone Learn This Stuff?}

\lettrine{E}{very book has its origin story}, and this one is no exception. If I were to capture the essence of creating this book in a single word, that word would be \textbf{curiosity} though \emph{improvised} comes as a close second. What you hold in your hands (or view on your screen) is the result of years of persistent questioning, a journey that began with a simple yet profound realization: 

\begin{center}
    \emph{What is an array?}
\end{center}
I didn't understand arrays. Not really.\\
I knew how to use them. I'd written plenty of code that used them. But if someone asked me what an array actually is not the syntax, not the API, but the thing itself I couldn't give a good answer. That bothered me.\\
So I started digging. What I found was that the question doesn't stop at the language level. It goes through the runtime, through memory layout, through the hardware. Follow it far enough and you end up at transistors, and voltages, and the question of how a physical thing can represent information at all.\\
That's what this book is about. It's one question followed seriously.
\section*{The Name and Its Meaning}
Arliz is an acronym: Arrays, Reasoning, Logic, Identity, Zero. It started as a name that sounded right, and the acronym came later. You can say it "ar-liz." It doesn't matter much. I personally say "ar-liz," but the pronunciation matters less than the journey it represents.

\vfill
{\small \linespread{1.2}\selectfont
\section*{One Work, Three Volumes}
\textit{Arliz} was never meant to be read in a single sitting, and after enough rewrites I stopped pretending otherwise. In fact, it was never even meant to become a book.\\
What began as a single article slowly grew as new ideas, revisions, and discoveries accumulated over the course of two years. The original manuscript expanded from a few pages into a much longer essay; the essay evolved into a multi-chapter book; the book itself eventually outgrew a single volume and was divided into several parts. Today, \textit{Arliz} spans three volumes:

\begin{itemize}
    \item \textbf{Volume I --- Zero to Bit.} How information gets encoded. Voltages, bits, integers, floats, characters, pointers, bytes.
    \item \textbf{Volume II --- Silicon Horizon.} The hardware that runs the code. Logic gates, memory, caches, CPUs, GPUs.
    \item \textbf{Volume III --- Array Odyssey.} Arrays themselves. Layout, algorithms, parallel systems, applications.
\end{itemize}
Each volume stands on its own, but they build on each other. Volume II assumes Volume I. Volume III assumes both.

\section*{What This Work Represents}
Arliz is not a gentle introduction to programming, nor is it a purely theoretical treatment of data structures. It is a comprehensive technical exploration of the most fundamental data structure in computing, examined from every relevant angle, and grounded the entire way down in the physical reality of the machines that run our code.\\
This is a living work. As I discover better explanations, uncover new implementation details, or recognize deeper connections between concepts, each volume improves sometimes independently of the others. Your engagement, through corrections, suggestions, and questions, makes you part of that evolution. The Introduction that follows is the practical companion to this more personal note: it covers prerequisites, difficulty, and how this volume fits into the larger work.

\vfill

\begin{flushright}
    \textit{Mahdi Mamashli} \\
    \textit{2026}
\end{flushright}
}
\clearpage